%!TEX root = ../main.tex

\chapter{Πειραματική Αξιολόγηση και Συμπεράσματα}   

\section{Δοκιμές Δυνατοτήτων του StarPUPy και torcpy}

Πριν προχωρήσουμε στην παρουσίαση των αποτελεσμάτων των πειραμάτων, και στα συμπεράσματα που προκύπτουν από αυτά, να πούμε ότι κάναμε και κάποιες δοκιμές για να δούμε κάποιες έξτρα δυνατότητες των δύο μοντέλων παράλληλου προγραμματισμού.

\subsection{StarPUPy: MPI Server-Client}

Σε αυτό το σημείο να πούμε ότι θέλαμε να δοκιμάσουμε και τις δυνατότητες του StarPUPy για εκτέλεση με MPI launcher, μέσω του Server-Client (ή Master-Slave σε παλαιότερες εκδόσεις) μοντέλου, για να είναι πιο κοντά η σύγκριση με το torcpy. Αλλά δυστυχώς η εκτέλεση αυτή των εφαρμογών δεν ήταν εφικτή λόγω προβλημάτων και έλλειψη πλήρης υποστήριξης των δυνατοτήτων MPI που παρέχει κανονικά το StarPU στην γλώσσα C. Το πρόβλημα ήταν κυρίως ότι δεν υποστηρίζει την χρήση native δομών δεδομένων της Python (π.χ. dictionaries) παρά μόνο NumPy arrays, και αυτό απαιτούσε να γίνει μετατροπή των δεδομένων σε NumPy arrays για να μπορέσει να γίνει η εκτέλεση μέσω MPI launcher. Έτσι, δοκιμάσαμε να μετατρέψουμε έστω την εφαρμογή Monte Carlo (app02) να χρησιμοποιεί NumPy arrays για να εκτελεστεί μέσω MPI launcher, αλλά η μετατροπή αυτή δεν ήταν αρκετή γιατί μετά προέκυψαν άλλα προβληματα με το Initialization του περιβάλλοντος MPI (Server-Client) σε συνδυασμό με το StarPU. Έτσι, καταλήξαμε ότι η δυνατότητες MPI δεν έχουν μεταφερθεί ακόμα πλήρως στην έκδοση της StarPU για Python και το αφήσαμε εκεί. Ο κώδικας για αυτή την έκδοση της εφαρμογής Monte Carlo (app02) υπάρχει στο αποθετήριο του έργου, αλλά δεν είναι πλήρως λειτουργικός και δεν χρησιμοποιήθηκε για την αξιολόγηση των μοντέλων παράλληλου προγραμματισμού \cite{appcode}. 

\subsection{Εμφολευμένη Παραλληλοποίηση (Nested Parallelism)}

Επιπλέον, δοκιμάσαμε να δούμε την δυνατότητα εμφολευμένης παραλληλοποίησης (nested parallelism) εργασιών των δύο μοντέλων και πώς αυτά το αντιμετωπίζουν. Έτσι φτιάξαμε μια νέα εκδοχή της εφαρμογής Image Processing (app01) που εκτελείται σε δύο επίπεδα παραλληλοποίησης. Στο πρώτο επίπεδο, έχουμε δύο εργασίες που η κάθε μία παίρνει τις μισές εικόνες του συνόλου των εικόνων και ως δεύτερο επίπεδο παραλληλοποίησης, κάθε εργασία δημιουργεί τις υπο-εργασίες που εκτελούν την ανάλυση υφής με φίλτρα Gabor για τις εικόνες που τους ανατίθενται. Έτσι, έχουμε δύο επίπεδα παραλληλοποίησης, όπου στο πρώτο επίπεδο γίνεται διαμοιρασμός των εικόνων σε δύο εργασίες και στο δεύτερο επίπεδο γίνεται διαμοιρασμός των εικόνων που ανατίθενται σε κάθε εργασία σε υπο-εργασίες. Ο κώδικας για αυτή την εκδοχή της εφαρμογής Image Processing (app01) υπάρχει στο αποθετήριο του έργου, αλλά δεν χρησιμοποιήθηκε για την αξιολόγηση των μοντέλων παράλληλου προγραμματισμού \cite{appcode}.

Από την εκτέλεση αυτής της εκδοχής της εφαρμογής Image Processing (app01) προέκυψε ότι
και τα δύο μοντέλα παράλληλου προγραμματισμού (StarPUPy και torcpy) υποστηρίζουν την εμφολευμένη παραλληλοποίηση, αλλά με μία μιρκή διαφορά. Το torcpy υποστηρίζει την εμφολευμένη παραλληλοποίηση απαυθείας χωρίς κάποια έξτρα ευθύνη συχρονισμού ή εξασφάλιση αρκούντων πόρων για την εκτέλεση των υπο-εργασιών, ενώ το StarPUPy απαιτεί από τον προγραμματιστή να εξασφαλίσει ότι υπάρχουν αρκετοί διαθέσιμοι πόροι για την εκτέλεση των υπο-εργασιών. Διότι στον κώδικα που έχουμε, έχουμε στην ουσία 2 μητρικές εργασίες, η κάθε μία δημιουργεί από 50 περιπού υπο-εργασίες όσες και οι είκονες που έχει. Αν εκτελέσουμε το πρόγραμμα με λιγότερο από 3 νήματα εργασίας, τότε το StarPUPy δεν θα μπορεί να εκτελέσει τις υπο-εργασίες και θα περιμένει για πάντα. Αντίθετα, το torcpy θα εκτελέσει τις υπο-εργασίες με τα διαθέσιμα νήματα και θα ολοκληρωθεί κανονικά. Αυτό γίνεται διότι μετά το submit των υπο-εργασιών, το torcpy δεν περιμένει να ολοκληρωθούν οι υπο-εργασίες για να συνεχίσει (non-blocking wait), ενώ το StarPUPy περιμένει να ολοκληρωθούν οι υπο-εργασίες πριν συνεχίσει (blocking wait). Έτσι, το StarPUPy απαιτεί από τον προγραμματιστή να εξασφαλίσει ότι υπάρχουν αρκετοί διαθέσιμοι πόροι για την εκτέλεση των υπο-εργασιών, ενώ το torcpy δεν απαιτεί κάτι τέτοιο.

\section{Αποτελέσματα Πειραμάτων και Συμπεράσματα}

Αρχικά, στο Σχήμα \ref{fig:app_results_standard} παρουσιάζονται οι χρόνοι εκτέλεσης και οι χρονοβελτίωση όλων των εφαρμογών που αναπτύξαμε, όταν αυτές εκτελέστηκαν με το Standard (NUMA-oblivious) τρόπο εκτέλεσης.

\begin{figure}[H]
    \centering
    \includegraphics[width=\textwidth]{figures/fig1_framework_comparison_standard}
    \caption{Αποτελέσματα Πειραμάτων - Standard (NUMA-oblivious) τρόπος εκτέλεσης.}
    \label{fig:app_results_standard}
\end{figure}

Τα πρώτα τρία γραφήματα του παραπάνω σχήματος παρουσιάζουν τους χρόνους εκτέλεσης των εφαρμογών. Ο άξονας y αντιπροσωπεύει τον χρόνο εκτέλεσης σε δευτερόλεπτα, ενώ ο άξονας x αντιπροσωπεύει τον αριθμό των νήματων εργασίας που χρησιμοποιήθηκαν για την εκτέλεση των εφαρμογών. Αλλά επειδή για το torcpy έχουμε συνδυασμούς εκτέλεσης με διεργασίες και νήματα αντί καθαρά νήματα όπως στο StarPUPy, πήραμα από κάθε συνδυασμό εκτέλεσης (cores=processes x threads) τον καλύτερο χρόνο. Δηλαδή, για το torcpy, αν εκτελέσαμε την εφαρμογή με 4 νήματα και 2 διεργασίες, και με 2 νήματα και 4 διεργασίες, τότε κρατήσαμε τον καλύτερο χρόνο από τους δύο συνδυασμούς. Έτσι, τα γραφήματα παρουσιάζουν τον καλύτερο χρόνο εκτέλεσης για κάθε αριθμό νήματων εργασίας που χρησιμοποιήθηκαν.

Στα από κάτω τρία γραφήματα του Σχήματος \ref{fig:app_results_standard} παρουσιάζεται η χρονοβελτίωση (speedup) των εφαρμογών για κάθε αριθμό νήματων εργασίας σε σχέση με την σειριακή εκτέλεση (1 process x 1 thread). Ο άξονας y αντιπροσωπεύει την χρονοβελτίωση σε ποσοστό, ενώ ο άξονας x αντιπροσωπεύει τον αριθμό των νήματων εργασίας που χρησιμοποιήθηκαν για την εκτέλεση των εφαρμογών, όπως και στα από πάνω γραφήματα. Επιπλέον, σε κάθε γράφημα υπάρχει και η θεωρητική ιδανική χρονοβελτίωση (ideal speedup) που θα είχαμε αν η εκτέλεση ήταν τέλεια παραλληλοποιημένη και δεν είχαμε κανένα είδους καθυστέρηση. 

Πριν προχωρήσουμε στην ανάλυση των αποτελεσμάτων, έχουμε και το Σχήμα \ref{fig:parallel_efficiency} όπου παρουσιάζεται η παράμετρος της παραλληλικής αποδοτικότητας (parallel efficiency) για κάθε εφαρμογή και για κάθε αριθμό νήματων εργασίας που χρησιμοποιήθηκαν. Η παράμετρος της παραλληλικής αποδοτικότητας υπολογίζεται ως ο λόγος της χρονοβελτίωσης προς τον αριθμό των νήματων εργασίας που χρησιμοποιήθηκαν. Δηλαδή, αν έχουμε χρονοβελτίωση 4 με 8 νήματα εργασίας, τότε η παραλληλική αποδοτικότητα είναι 4/8 = 0.5 ή 50\%. Ουσιαστικά, η παραλληλική αποδοτικότητα μας δείχνει πόσο καλά εκμεταλλευόμαστε τους διαθέσιμους πόρους (νήματα εργασίας) για την εκτέλεση των εφαρμογών.

\begin{figure}[H]
    \centering
    \includegraphics[width=\textwidth]{figures/fig4_efficiency_comparison}
    \caption{Παράμετρος Παραλληλικής Αποδοτικότητας (Parallel Efficiency) για κάθε εφαρμογή και για κάθε αριθμό νήματων εργασίας.}
    \label{fig:parallel_efficiency}
\end{figure}

\subsection{Ανάλυση Αποτελεσμάτων Πειραμάτων}

\subsubsection{CMA-ES (app00)}

Από το Σχήμα \ref{fig:app_results_standard} παρατηρούμε ότι το StarPUPy υπερέχει σημαντικά του torcpy, επιτυγχάνοντας ελάχιστο χρόνο εκτέλεσης στους 16 πυρήνες, ενώ το torcpy εμφανίζει χαμηλότερο speedup (περίπου 3.5x) και αυξάνει τον χρόνο εκτέλεσης μετά τους 16 πυρήνες. Η παράλληλη αποδοτικότητα (Σχήμα \ref{fig:parallel_efficiency}) του torcpy πέφτει κατακόρυφα μετά τους 4 πυρήνες, ενώ του StarPUPy παραμένει πάνω από 80\% μέχρι τους 16 επεξεργαστές.

Αυτό προκύπτει από το γεγονός ότι η εφαρμογή αυτή είναι sync-bound (περιορίζεται από συγχρονισμό). Όπως φαίνεται στον κώδικα, η παραλληλοποίηση γίνεται αυστηρά εντός κάθε γενιάς (γενιά $g$ προς $g+1$) με έναν βρόχο \texttt{while not es.stop()}. Το StarPUPy χρησιμοποιεί ελαφριά, ασύγχρονα νήματα (C-backend) και asyncio.futures που επιτρέπουν γρήγορο barrier συγχρονισμό ανά γενιά. Ενώ, το torcpy χρησιμοποιεί τη ρουτίνα \texttt{torc.map()}, η οποία σε κάθε γενιά εισάγει overhead κατανομής εργασιών μέσω διεργασιών MPI. Καθώς αυξάνονται οι πυρήνες, το κόστος συγχρονισμού του MPI barrier και το GIL thrashing για την ανάκτηση των αποτελεσμάτων ανά γενιά κυριαρχούν έναντι του μικρού υπολογιστικού έργου (πλέγμα $128 \times 128$) \cite{appcode}.

\subsubsection{Image Processing - Gabor Filters (app01)}

Στην δεύτερη εφαρμογή, Επεξεργασία Εικόνας - Φίλτρα Gabor (app01), παρατηρούμε ότι γίνεται το αντίστροφο σενάριο. Το torcpy εμφανίζει εξαιρετική κλιμάκωση (speedup ~22x στους 64 πυρήνες) και η αποδοτικότητά του μειώνεται ομαλά. Το StarPUPy καταρρέει πλήρως μετά τους 8 πυρήνες, με τον χρόνο εκτέλεσης να αυξάνεται εκθετικά και την αποδοτικότητα να πέφτει κοντά στο 0\%. 

Εδώ βασικό ρόλο παίζει πάλι η φύση της εφαρμογής. Η εφαρμογή αυτή είναι memory-bound (περιορίζεται από το εύρος ζώνης μνήμης) και εκτελεί 40 FFT convolutions ανά εικόνα. Το StarPUPy εκτελείται αποκλειστικά με threads εντός του ίδιου Python interpreter. Παρά το γεγονός ότι οι βιβλιοθήκες FFT (NumPy/BLAS) απελευθερώνουν το GIL, η ταυτόχρονη πρόσβαση δεκάδων νημάτων στην ίδια κοινή μνήμη προκαλεί ακραίο contention. Ενώ, το torcpy υπερέχει επειδή βασίζεται σε διεργασίες MPI. Κάθε worker thread ανήκει σε απομονωμένη διεργασία με δικό της χώρο διευθύνσεων (address space) αποφεύγοντας μειώντας έτσι το contention.

\subsubsection{Monte Carlo (app02)}

Και τα δύο συστήματα εμφανίζουν παρόμοια, ικανοποιητική κλιμάκωση μέχρι τους 8-16 πυρήνες. Μετά τους 16 πυρήνες, το torcpy καταφέρνει να διατηρήσει μια ελαφρώς καλύτερη καμπύλη speedup, ενώ το StarPUPy παρουσιάζει φαινόμενα κορεσμού (saturation).

Η εφαρμογή είναι compute-bound (έντονοι μαθηματικοί υπολογισμοί με Floops: exp, sin, log1p). Όπως φαίνεται στον κώδικα, η επικοινωνία είναι ελάχιστη, καθώς κάθε batch επιστρέφει μόνο τρεις scalar τιμές \texttt{(out\_sum, out\_sumsq, out\_n)}, πραγματοποιώντας ένα πολύ αποδοτικό parallel reduction στο τέλος. Η μικρή υπεροχή του torcpy στους πολλούς πυρήνες οφείλεται στην πλήρη εξάλειψη της πιθανότητας GIL thrashing λόγω των διακριτών MPI διεργασιών \cite{appcode}.

\subsection{NUMA-aware εκτέλεση των εφαρμογών}

